\section{局部环与切空间}\label{sec:local-rings-and-tangent-spaces}

现在把切空间和光滑性的概念从仿射情形推广到拟射影情形. 初次尝试定义拟射影簇在一点的切空间时, 可能只是限制到该点的某个仿射邻域, 然后使用对应仿射簇在该点的切空间. 这么做确实给出正确的概念, 但它有一个不足: 不同的仿射限制是否给出同一个 (典范同构的) 切空间, 并不立即清楚. 为了避开这一微妙之处, 我们使用局部环建立研究切空间的另一种方法.

大致说来, 拟射影簇 $X$ 在点 $a\in X$ 处的局部环, 表现得像 $a$ 的一个无穷小邻域的坐标环, 这个邻域包含在所有其他邻域中. 局部环是刻画 $X$ 在 $a$ 附近的几何信息的代数工具. 特别地, 可以用它以内在方式定义 $X$ 在 $a$ 处的切空间, 自然地推广 Theorem~\ref{thm:point-ideal-square-quotient-is-dual-of-tangent-space} 中切空间的内在描述. 如将看到的, 这给出了拟射影簇在一点的切空间的一个定义, 当限制到任一仿射邻域时, 它恢复仿射情形中的概念.

我们从一点处局部环的正式定义开始讨论.

\begin{Definition}{拟射影簇在一点处的局部环}\leavevmode\label{def:local-ring-of-quasiprojective-variety-at-point}
令 $X$ 为拟射影簇, $a\in X$. $X$ 在 $a$ 处的\emph{局部环}记为 $K[X]_{a}$, 是数对 $(U,F)$ 的等价类的集合, 其中 $U\subseteq X$ 是包含 $a$ 的 $X$ 的开子集, $F:U\to K$ 是正则函数. 两个数对 $(U,F)$ 与 $(V,G)$ 等价, 当且仅当存在包含 $a$ 的开集 $W\subseteq U\cap V$, 使对所有 $b\in W$ 有 $F(b)=G(b)$.
\end{Definition}

换句话说, $K[X]_{a}$ 的一个元素可由 $a$ 的某个邻域上的正则函数表示. 两个可能定义在不同邻域上的正则函数给出同一个元素, 当且仅当它们在某个可能更小的邻域上相等. 因此, 可把 $X$ 在 $a$ 处的局部环的一个元素看成定义在 $a$ 的任意小邻域上的正则函数. 为此, 把局部环 $K[X]_{a}$ 的元素称为 $a$ 处的\emph{局部函数}.

Definition~\ref{def:local-ring-of-quasiprojective-variety-at-point} 有几个方面需要进一步说明. 例如, 应当停下来理解为什么数对 $(U,F)$ 之间的关系确实是等价关系, 也应仔细定义 $K[X]_{a}$ 上自然的加法, 乘法和标量乘法, 这些代数运算使它成为 $K$-代数 (见 Exercise~\ref{ex:local-ring-equivalence-classes-and-algebra-operations}).

在进一步研究局部环前, 先考察几个具体例子.

\begin{Example}{由多项式的商给出的局部函数}\leavevmode\label{ex:local-functions-from-homogeneous-quotients}
令 $X\subseteq\mathbb{P}^{n}$ 为拟射影簇, $a\in X$. 给定同次数的两个齐次多项式 $f,g\in K[x_{0},\ldots,x_{n}]$, 它们的商定义正则映射 $f/g:X\setminus\mathcal{V}(g)\to K$. 特别地, 若 $g(a)\ne0$, 则 $f/g$ 给出 $a$ 处的一个局部函数, 它定义在邻域 $U=X\setminus\mathcal{V}(g)$ 上. 事实上, 下面将看到每个局部函数都可以实现为多项式的商.
\end{Example}

\begin{Example}{不同的商给出同一个局部函数}\leavevmode\label{ex:local-function-does-not-see-other-irreducible-component}
考察可约射影簇 $X=\mathcal{V}(xy)=\mathcal{V}(x)\cup\mathcal{V}(y)\subseteq\mathbb{P}^{2}$, 令 $a=[1:0:0]$. 注意 $a$ 在不可约分支 $\mathcal{V}(y)$ 中, 但不在不可约分支 $\mathcal{V}(x)$ 中. 考察有理函数
\[
  \frac{x+y}{x-y}:X\setminus\{[0:0:1]\}\to K.
\]
这个有理函数并非恒定, 因为它在 $[1:0:0]$ 处取值 $1$, 在 $[0:1:0]$ 处取值 $-1$. 不过, 它定义的 $a$ 处的局部函数是常值函数 $1$, 原因只在于限制到邻域 $U=X\setminus\mathcal{V}(x)$ 后, 它的分子和分母相同. 这一例子说明局部函数的局部性质: $a$ 处的局部函数\textquotedbl{}看不见\textquotedbl{}不可约分支 $\mathcal{V}(x)\subseteq X$, 因为这个不可约分支不含 $a$.
\end{Example}

以先前例子中的观察为基础, 下一个结果说明拟射影簇 $X$ 在点 $a$ 处的局部环确实只取决于 $a\in X$ 附近的情况: 换句话说, 将其限制到 $a$ 的更小邻域并不影响它.

\begin{Proposition}{局部环只取决于局部}\leavevmode\label{prop:local-ring-unchanged-on-open-neighborhood}
设 $X$ 为拟射影簇, $a\in X$. 若 $U\subseteq X$ 为包含 $a$ 的开子簇, 则
\[
  K[U]_{a}=K[X]_{a}.
\]
\end{Proposition}
\begin{Proof}
令 $X$ 为拟射影簇, $a\in X$, 令 $U\subseteq X$ 为包含 $a$ 的开子簇. 有典范的局部环限制同态
\[
  K[X]_{a}\to K[U]_{a},
\]
它将局部函数 $(V,F)\in K[X]_{a}$ 送到 $(U\cap V,F|_{U\cap V})\in K[U]_{a}$. 注意, 局部函数为 $0$, 当且仅当代表它的任意正则函数在 $a$ 的某个邻域上消没; 这意味着限制映射是单射. 限制映射也是满射, 因为按定义, $K[U]_{a}$ 的每个元素都可视为 $K[X]_{a}$ 的元素, 因而是它自身的限制. 因此 $K[X]_{a}=K[U]_{a}$.
\end{Proof}

由于拟射影簇是局部仿射的, Proposition~\ref{prop:local-ring-unchanged-on-open-neighborhood} 把局部环的研究归结到仿射情形. 因此, 具体描述仿射簇的局部环将是有用的, 现在给出这一描述.

\begin{Lemma}{仿射簇的局部环}\leavevmode\label{lem:local-rings-of-affine-varieties-as-fractions}
令 $X$ 为仿射簇, $a\in X$. $K[X]_{a}$ 的每个元素都可表示为商 $F/G$, 其中 $F,G\in K[X]$ 为正则函数, 且 $G(a)\ne0$. 两个商 $F/G$ 和 $F'/G'$ 给出同一个局部函数, 当且仅当存在 $H\in K[X]$, 使 $H(a)\ne0$, 且
\[
  H(FG'-F'G)=0\in K[X].
\]
\end{Lemma}
\begin{Proof}
由于开集以及开集上的正则函数环在同构下保持, 只须假设 $X$ 为 $\mathbb{A}^{n}_{0}\subseteq\mathbb{P}^{n}$ 的闭子集. 令 $a\in X$, 假设 $(U,H)$ 代表 $K[X]_{a}$ 的一个元素, 因而 $H\in K[U]$. 由正则函数的定义, 存在同为 $d$ 次的齐次多项式 $f,g\in K[x_{0},\ldots,x_{n}]$, 使 $g(a)\ne0$, 且
\[
  H(b)=\frac{f(b)}{g(b)}\qquad\text{对所有 }b\in U\setminus\mathcal{V}(g).
\]
定义 $F=f(x_{0},x_{1},\ldots,x_{n})/x_{0}^{d}$ 和 $G=g(x_{0},\ldots,x_{n})/x_{0}^{d}$. 由于 $F,G$ 各自是同次数的齐次多项式的商, 且分母在 $X$ 上不消没, 可知 $F,G\in K[X]$. 由构造, 有
\[
  H(b)=\frac{F(b)}{G(b)}\qquad\text{对所有 }b\in U\setminus\mathcal{V}(G).
\]
这表明在 $K[X]_{a}$ 中, $(U,H)$ 与 $(X\setminus\mathcal{V}(G),F/G)$ 等价, 所以每个局部函数都可实现为 $X$ 上两个正则函数的商.

现假设 $F/G$ 和 $F'/G'$ 是 $X$ 上正则函数的商, 在 $K[X]_{a}$ 中给出同一个局部函数. 则 $FG'-F'G$ 必在 $a$ 的某个邻域上消没. 由于不可约分支的开子集是稠密的, 且 $FG'-F'G$ 的消没集在 $X$ 中闭, 这蕴含 $FG'-F'G$ 在 $X$ 的每个包含 $a$ 的不可约分支上消没. 令 $H\in K[X]$ 为任一满足 $H(a)\ne0$ 且在 $X$ 的每个不包含 $a$ 的不可约分支上消没的正则函数. 则 $H(FG'-F'G)=0\in K[X]$.

反过来, 假设存在 $F,F',G,G',H\in K[X]$, 后三者均不在 $a$ 处消没, 且 $H(FG'-F'G)=0\in K[X]$. 因而
\[
  F(b)/G(b)=F'(b)/G'(b)\qquad\text{对所有 }b\in X\setminus\mathcal{V}(HGG').
\]
由于 $X\setminus\mathcal{V}(HGG')$ 是 $a$ 的开邻域, 可知 $F/G$ 和 $F'/G'$ 在 $K[X]_{a}$ 中给出同一个局部函数.
\end{Proof}

\BookPairedNoteLeft{
Lemma~\ref{lem:local-rings-of-affine-varieties-as-fractions} 中的描述是更一般的代数过程\textquotedbl{}局部化\textquotedbl{}的动机; 更多细节见 Exercise~\ref{ex:construction-and-examples-of-ring-localization}.
}{
现在从局部环的角度重新考察仿射簇的切空间. 回顾第~\ref{chap:smoothness}~章, 仿射簇 $X$ 在点 $a\in X$ 处的切空间 $T_{a}X$, 首先用切向量和 $a$ 处的梯度作外在定义, 后来又用 $X$ 上在 $a$ 处消没的正则函数构成的极大理想作内在描述.
}
第二种内在观点在局部环的情形有自然的对应, 下面说明.

\BookPairedNote{
令 $X$ 为拟射影簇, $a\in X$, 令 $J_{a}\subseteq K[X]_{a}$ 为在 $a$ 处消没的局部函数构成的理想. 注意, $J_{a}$ 是 $K[X]_{a}$ 的极大理想. 事实上, 它是 $K[X]_{a}$ 中唯一的极大理想 (Exercise~\ref{ex:unique-maximal-ideal-of-local-ring}).
}{
受这一几何情形的启发, 具有唯一极大理想的环通常称为\textquotedbl{}局部环\textquotedbl{}.
}
专门来看仿射情形, 第~\ref{chap:smoothness}~章已经定义了切空间, 现在可用理想 $J_{a}\subseteq K[X]_{a}$ 对切空间作如下另一种描述.

\begin{Proposition}{再论仿射簇的切空间}\leavevmode\label{prop:affine-tangent-space-dual-of-local-maximal-ideal-quotient}
若 $X$ 为仿射簇, $a\in X$, $J_{a}\subseteq K[X]_{a}$ 为在 $a$ 处消没的局部函数的理想, 则有典范的向量空间同构
\[
  T_{a}X=(J_{a}/J_{a}^{2})^{\vee}.
\]
\end{Proposition}
\begin{Proof}
如第~\ref{sec:intrinsic-tangent-spaces}~节, 令 $I_{a}\subseteq K[X]$ 表示在 $a$ 处消没的正则 (即多项式) 函数的极大理想. 由 Theorem~\ref{thm:point-ideal-square-quotient-is-dual-of-tangent-space}, 有典范的向量空间同构
\[
  T_{a}X=(I_{a}/I_{a}^{2})^{\vee}.
\]
因此, 我们必须证明存在典范的向量空间同构
\[
  I_{a}/I_{a}^{2}=J_{a}/J_{a}^{2}.
\]
考察限制同态 $K[X]\to K[X]_{a}$. 稍加思考可知此映射把 $I_{a}$ 送入 $J_{a}$, 把 $I_{a}^{2}$ 送入 $J_{a}^{2}$. 因而得到典范的线性映射
\begin{equation}\label{eq:restriction-map-of-cotangent-quotients}
  I_{a}/I_{a}^{2}\to J_{a}/J_{a}^{2},
\end{equation}
只须再证明这个线性映射是双射.

为证明 \eqref{eq:restriction-map-of-cotangent-quotients} 的满射性, 首先注意, 由 Lemma~\ref{lem:local-rings-of-affine-varieties-as-fractions}, $K[X]_{a}$ 的每个元素都可写成 $F/G$, 其中 $F,G\in K[X]$ 且 $G(a)\ne0$. 假设 $F/G$ 给出 $J_{a}$ 的一个元素, 或等价地, $F\in I_{a}$. 必要时同时缩放 $F$ 和 $G$, 可假设 $G(a)=1$. 因而 $F$ 与 $1/G-1$ 都给出 $J_{a}$ 的元素, 所以
\[
  F/G-F=F(1/G-1)\in J_{a}^{2}.
\]
这意味着 $F/G$ 与 $F$ 定义的局部函数在商 $J_{a}/J_{a}^{2}$ 中等价, 所以 $[F/G]\in J_{a}/J_{a}^{2}$ 是 $[F]\in I_{a}/I_{a}^{2}$ 的像, 从而证得 \eqref{eq:restriction-map-of-cotangent-quotients} 的满射性.

为证明单射性, 假设 $F\in I_{a}$ 限制为属于 $J_{a}^{2}$ 的一个局部函数. $J_{a}^{2}$ 的元素是 $J_{a}$ 中两个元素的乘积的有限和. 因此由 Lemma~\ref{lem:local-rings-of-affine-varieties-as-fractions}, 存在 $m\ge1$ 和 $i=1,\ldots,m$ 时的 $G_{i},G'_{i},H_{i},H'_{i}\in K[X]$, 使对所有 $i=1,\ldots,m$ 有 $G_{i}(a)=0=G'_{i}(a)$ 且 $H_{i}(a)\ne0\ne H'_{i}(a)$, 并存在 $a$ 的邻域 $U$, 使
\[
  F(b)=\sum_{i=1}^{m}\frac{G_{i}(b)G'_{i}(b)}{H_{i}(b)H'_{i}(b)}
  \qquad\text{对所有 }b\in U.
\]
选取 $H\in K[X]$, 使 $H(a)\ne0$ 且 $H$ 在 $X\setminus U$ 上消没. 则
\[
  H(b)F(b)=H(b)\sum_{i=1}^{m}\frac{G_{i}(b)G'_{i}(b)}{H_{i}(b)H'_{i}(b)}
  \qquad\text{对所有 }b\in X.
\]
消去分母后可见 $(H_{1}\cdots H_{m})(H'_{1}\cdots H'_{m})HF\in I_{a}^{2}$. 由于 $F$ 是此乘积中唯一属于 $I_{a}$ 的因子, 可推出 $F\in I_{a}^{2}$ (Exercise~\ref{ex:divide-out-factor-not-in-maximal-ideal-from-square}). 这证明了 \eqref{eq:restriction-map-of-cotangent-quotients} 的单射性, 从而完成证明.
\end{Proof}

前一结果启发了拟射影簇一般情形下切空间的如下内在定义.

\begin{Definition}{拟射影簇的切空间}\leavevmode\label{def:tangent-space-of-quasiprojective-variety}
设 $X$ 为拟射影簇, $a\in X$, 令 $J_{a}$ 表示 $K[X]_{a}$ 中在 $a$ 处消没的局部函数构成的理想. $X$ 在 $a$ 处的\emph{切空间}记为 $T_{a}X$, 是向量空间
\[
  T_{a}X=(J_{a}/J_{a}^{2})^{\vee}.
\]
\end{Definition}

由 Definition~\ref{def:tangent-space-of-quasiprojective-variety} 的局部性, 若 $X$ 为拟射影簇, $a\in X$, $U$ 为 $a$ 的一个邻域, 则 $T_{a}X=T_{a}U$. 特别地, 总可以先限制到一个仿射开邻域, 然后按第~\ref{chap:smoothness}~章中的方法计算拟射影簇的切空间.

由切空间在仿射情形的内在性质还可知, 切空间在拟射影情形也是内在的. 更具体地说, 若 $F:X\to Y$ 是拟射影簇的同构, 且 $F(a)=b$, 则 $F$ 限制为 $a$ 的任一仿射邻域与其在 $Y$ 中的像之间的同构, 而这个像是 $b$ 的仿射邻域. 由仿射情形中切空间的内在性质, 可知 $T_{a}X=T_{b}Y$.

有了拟射影情形中的切空间概念, 自然可以比较切空间的向量空间维数与簇的维数. 通过归结到仿射情形, 首先可观察到, 只要 $X$ 不可约, 就有
\[
  \dim(T_{a}X)\ge\dim(X)
\]
(Exercise~\ref{ex:quasiprojective-tangent-dimension-at-least-variety-dimension}). 与仿射情形一样, 等号成立的点尤其重要.

\begin{Definition}{拟射影簇的光滑性}\leavevmode\label{def:smoothness-of-quasiprojective-variety}
令 $X$ 为不可约拟射影簇, $a\in X$. 若 $\dim(T_{a}X)=\dim(X)$, 就称 $X$ 在 $a$ 处\emph{光滑}; 否则称 $X$ 在 $a$ 处\emph{奇异}. 若 $X$ 在每个点 $a\in X$ 处均光滑, 则称 $X$ \emph{光滑}; 否则称 $X$ \emph{奇异}.
\end{Definition}

由于我们在一般拟射影情形下对切空间和光滑性的建立是局部的, 对这些概念的直观理解可以完全来自仿射情形. 特别地, 拟射影簇在一点处光滑还是奇异, 可以通过限制到该点的一个仿射邻域判定. 因而仿射情形中建立的结果, 例如光滑性的 Jacobian 判据和奇异点集为闭集的事实, 都自然地推广到拟射影情形. 请读者在习题中探究这些结果的射影对应.

注意, 射影情形下的 Jacobian 判据通常比预期多涉及一个偏导数. 本节最后讨论关于 $\mathbb{P}^{n}$ 中超曲面的一个 Jacobian 判据的特例, 由此可见齐次性如何帮助我们处理这个额外的偏导数.

\begin{Example}{射影超曲面的奇点}\leavevmode\label{ex:singularities-of-projective-hypersurfaces}
令 $f\in K[x_{0},\ldots,x_{n}]$ 为不可约齐次多项式, 考察射影超曲面 $X=\mathcal{V}(f)\subseteq\mathbb{P}^{n}$. 我们证明 $X$ 在点 $a\in X$ 处奇异, 当且仅当 $f$ 的所有偏导数都在 $a$ 处消没:
\[
  \frac{\partial f}{\partial x_{0}}(a)=\cdots=\frac{\partial f}{\partial x_{n}}(a)=0.
\]
不妨设 $a\notin\mathcal{V}(x_{0})$, 因而 $a=[1:a_{1}:\cdots:a_{n}]$. 限制到仿射邻域 $X_{0}=X\cap\mathbb{A}^{n}_{0}$, 作自然等同
\[
  X_{0}=\mathcal{V}(f_{0})\subseteq\mathbb{A}^{n}
  \qquad\text{其中}\qquad
  f_{0}(x_{1},\ldots,x_{n})=f(1,x_{1},\ldots,x_{n}).
\]
置 $a=(a_{1},\ldots,a_{n})\in\mathbb{A}^{n}$. 由于 $f$ 不可约, 从而 $f_{0}$ 不可约, 有 $\mathcal{I}(X_{0})=\langle f_{0}\rangle$. 根据仿射情形的 Jacobian 判据 (见 Example~\ref{ex:jacobian-criterion-for-hypersurfaces}), $X_{0}$ 在 $a$ 处奇异, 从而 $X$ 在 $a$ 处奇异, 当且仅当
\[
  \frac{\partial f_{0}}{\partial x_{1}}(a)=\cdots=\frac{\partial f_{0}}{\partial x_{n}}(a)=0.
\]
对偏导数稍作思考可知
\[
  \frac{\partial f}{\partial x_{i}}(1,a_{1},\ldots,a_{n})
  =\frac{\partial f_{0}}{\partial x_{i}}(a_{1},\ldots,a_{n})
  \qquad\text{对所有 }i=1,\ldots,n,
\]
因此 $X$ 在 $a\in X$ 处奇异, 当且仅当
\[
  \frac{\partial f}{\partial x_{1}}(a)=\cdots=\frac{\partial f}{\partial x_{n}}(a)=0.
\]
那么关于 $x_{0}$ 的偏导数呢? 由于 $f$ 是非常数齐次多项式, 齐次多项式的 Euler 恒等式 (Exercise~\ref{ex:euler-identity-for-homogeneous-polynomials}) 给出
\[
  \deg(f)f=x_{0}\frac{\partial f}{\partial x_{0}}+\cdots+x_{n}\frac{\partial f}{\partial x_{n}},
\]
又因为 $f(a)=0$ 且 $a\notin\mathcal{V}(x_{0})$, 可知所有偏导数在 $a$ 处消没, 当且仅当除 $\partial f/\partial x_{0}$ 外的偏导数都在 $a$ 处消没. 因此, $X$ 在 $a\in X$ 处奇异, 当且仅当 $f$ 的所有偏导数都在 $a\in X$ 处消没.
\end{Example}

\BookExerciseHeading
\begin{BookExercises}
  \item\label{ex:local-ring-equivalence-classes-and-algebra-operations}
  \begin{enumerate}
    \item 证明 Definition~\ref{def:local-ring-of-quasiprojective-variety-at-point} 中数对 $(U,F)$ 之间的关系是等价关系.
    \item 仔细定义局部函数的和, 积及标量倍数.
    \item 验证局部环 $K[X]_{a}$ 是 $K$-代数.
  \end{enumerate}
  \item\label{ex:unique-maximal-ideal-of-local-ring} 令 $X$ 为拟射影簇, $a\in X$. 证明 $K[X]_{a}$ 有唯一极大理想 $J_{a}\subseteq K[X]_{a}$, 由所有在 $a$ 处消没的局部函数构成.
  \item\label{ex:divide-out-factor-not-in-maximal-ideal-from-square} 令 $I$ 为环 $R$ 的极大理想, $a,b\in R$ 满足 $a\in I$, $b\notin I$, 且 $ab\in I^{2}$. 证明 $a\in I^{2}$.
  \item\label{ex:quasiprojective-tangent-dimension-at-least-variety-dimension} 令 $X$ 为不可约拟射影簇, $a\in X$. 证明
  \[
    \dim(T_{a}X)\ge\dim(X).
  \]
  (提示: 归结到第~\ref{chap:smoothness}~章中的适当结果.)
  \item\label{ex:smooth-projective-conics-and-singular-projective-cubics}
  \begin{enumerate}
    \item 令 $f\in K[x,y,z]$ 为不可约的二次齐次多项式. 证明 $\mathcal{V}(f)\subseteq\mathbb{P}^{2}$ 光滑.
    \item 举例说明 (a) 不能推广到三次.
  \end{enumerate}
  \item\label{ex:euler-identity-for-homogeneous-polynomials} 令 $f\in K[x_{0},\ldots,x_{n}]$ 为非常数齐次多项式. 证明
  \[
    \deg(f)f=x_{0}\frac{\partial f}{\partial x_{0}}+\cdots+x_{n}\frac{\partial f}{\partial x_{n}}.
  \]
  \item\label{ex:projective-jacobian-criterion} 设 $X\subseteq\mathbb{P}^{n}$ 为不可约射影簇, 消没理想为 $\mathcal{I}(X)=\langle f_{1},\ldots,f_{m}\rangle$. 证明 $X$ 在 $a\in X$ 处奇异, 当且仅当
  \[
    \operatorname{rk}(\operatorname{Jac}_{f_{1},\ldots,f_{m}}(a))<\operatorname{codim}(X).
  \]
  (注: $\operatorname{Jac}_{f_{1},\ldots,f_{m}}(a)$ 不是良定义的, 因为缩放 $a$ 的坐标会缩放 Jacobian 矩阵的各项; 但 Jacobian 的秩是良定义的, 因为这一缩放不影响秩.)
  \item\label{ex:singular-locus-of-irreducible-quasiprojective-variety-is-closed} 证明不可约拟射影簇的奇异点构成一个闭子簇.
  \item\label{ex:construction-and-examples-of-ring-localization} 令 $R$ 为环, $S\subseteq R$ 为对乘法封闭的子集.
  \begin{enumerate}
    \item 考察 $R\times S$ 上的如下关系: 若存在 $t\in S$ 使 $t(rs'-r's)=0$, 则 $(r,s)\sim(r',s')$. 证明 $\sim$ 是等价关系. 等价类的集合称为 $R$ 在 $S$ 处的\emph{局部化}, 记为 $R_{S}$.
    \item 证明 $R_{S}$ 是环, 其运算定义为
    \[
      \begin{aligned}
      (r,s)+(r',s')&=(rs'+r's,ss'),\\
      (r,s)(r',s')&=(rr',ss').
      \end{aligned}
    \]
    此外, 若 $R$ 为 $K$-代数, 证明 $R_{S}$ 为 $K$-代数, 其标量乘法为
    \[
      a(r,s)=(ar,s).
    \]
    \item 若 $R$ 为整环且 $S=R\setminus\{0\}$, 证明
    \[
      R_{S}=\operatorname{Frac}(R).
    \]
    \item 若 $R$ 为环, 且对某个非幂零元 $a\in R$, 有 $S=\{1,a,a^{2},\ldots\}$, 证明
    \[
      R_{S}=R[x]/\langle ax-1\rangle.
    \]
    \item 对某个仿射簇 $X$, 令 $R=K[X]$, 令 $I_{a}\subseteq R$ 为在 $a\in X$ 处消没的函数构成的极大理想. 若 $S=R\setminus I_{a}$, 证明
    \[
      R_{S}=K[X]_{a}.
    \]
  \end{enumerate}
\end{BookExercises}
